\begin{PSbox}[unbreakable]{obj}{اهداف یادگیری}

پس از تکمیل این ماژول، دانشجویان قادر خواهند بود:

\begin{enumerate}
\def\labelenumi{\arabic{enumi}.}
\tightlist
\item
  میان جامعه و نمونه تمایز قایل شوند
\item
  میان پارامتر و آماره تمایز قایل شوند
\item
  میانگین نمونه و واریانس نمونه را تعریف و محاسبه کنند
\item
  توزیع‌های نمونه‌گیری و اهمیت آن‌ها را توضیح دهند
\item
  گذار از احتمال \lr{$(\text{\rl{مدل}}\ \to \text{\rl{پیش‌بینی}})$} به آمار \lr{$(\text{\rl{داده}}\ \to \text{\rl{مدل}})$} را درک کنند
\end{enumerate}

\end{PSbox}

\PSsection{12.1}{مقدمه: از احتمال به آمار}

\PSsubsection{دو جهت}

{\def\LTcaptype{none} % do not increment counter
\begin{longtable}[c]{@{}
  >{\centering\arraybackslash}p{(0.965\linewidth - 2\tabcolsep) * \real{0.5217}}
  >{\centering\arraybackslash}p{(0.965\linewidth - 2\tabcolsep) * \real{0.4783}}@{}}
\toprule\noalign{}
\rowcolor{PSnavy}\PSth{احتمال} & \PSth{آمار} \\
\midrule\noalign{}
\endhead
\bottomrule\noalign{}
\endlastfoot
مدل \textbf{معلوم} است & مدل \textbf{نامعلوم} است \\
پیش‌بینی خروجی‌ها & استنباط مدل از داده \\
مسیله مستقیم & مسیله معکوس \\
«با داشتن توزیع، داده چگونه است؟» & «با داشتن داده، توزیع چگونه است؟» \\
\end{longtable}
}

\textbf{احتمال} (ماژول‌های \lr{\mbox{1–11}}): با داشتن یک مدل \lr{\mbox{\ensuremath{\to}}} مشاهده‌ها را پیش‌بینی کنید.\PSbr{}\textbf{آمار} (ماژول‌های \lr{\mbox{12–16}}): با داشتن مشاهده‌ها \lr{\mbox{\ensuremath{\to}}} مدل را استنباط کنید.

\PSsubsection{زمینه مهندسی}

\begin{itemize}
\tightlist
\item
  \textbf{احتمال:} «اگر نویز \lr{$N(0,\, 0.01)$} باشد، \lr{$P(\PSmtext{error})$} چقدر است؟»
\item
  \textbf{آمار:} «\lr{\mbox{1000}} نمونه نویز اندازه‌گیری کردم. \lr{$\sigma^{2}$} چقدر است؟»
\end{itemize}

\PSsection{12.2}{جامعه در برابر نمونه}

\PSsubsection{جامعه}

مجموعه \textbf{کامل} تمام اقلام/اندازه‌گیری‌های موردنظر.

{\def\LTcaptype{none} % do not increment counter
\begin{longtable}[c]{@{}
  >{\centering\arraybackslash}p{(0.965\linewidth - 2\tabcolsep) * \real{0.6129}}
  >{\centering\arraybackslash}p{(0.965\linewidth - 2\tabcolsep) * \real{0.3871}}@{}}
\toprule\noalign{}
\rowcolor{PSnavy}\PSth{زمینه مهندسی} & \PSth{جامعه} \\
\midrule\noalign{}
\endhead
\bottomrule\noalign{}
\endlastfoot
همه تراشه‌هایی که تاکنون در یک کارخانه تولید شده‌اند & \textbf{تمام} تأخیرهای تراشه \\
فرایند نویز تا ابد & \textbf{تمام} مقادیر ولتاژ نویز \\
تمام بسته‌های شبکه‌ای که تاکنون ارسال شده‌اند & \textbf{تمام} تأخیرهای بسته \\
\end{longtable}
}

جامعه‌ها معمولاً \textbf{نامتناهی یا به‌طور غیرعملی بزرگ} هستند \lr{—} نمی‌توانیم همه چیز را اندازه‌گیری کنیم.

\PSsubsection{نمونه}

یک \textbf{زیرمجموعه} از جامعه که در واقع مشاهده/اندازه‌گیری می‌کنیم.

{\def\LTcaptype{none} % do not increment counter
\begin{longtable}[c]{@{}cc@{}}
\toprule\noalign{}
\rowcolor{PSnavy}\PSth{زمینه مهندسی} & \PSth{نمونه} \\
\midrule\noalign{}
\endhead
\bottomrule\noalign{}
\endlastfoot
50 تراشه آزمایش‌شده از خط تولید & تأخیرهای اندازه‌گیری‌شده \\
10000 نمونه ولتاژ نویز ثبت‌شده & مقادیر ثبت‌شده \\
500 تأخیر بسته اندازه‌گیری‌شده & تأخیرهای مشاهده‌شده \\
\end{longtable}
}

\PSsubsection{چرا نمونه‌گیری می‌کنیم}

\begin{itemize}
\tightlist
\item
  نمی‌توان هر تراشه را آزمایش کرد (آزمون مخرب)
\item
  نمی‌توان نویز را تا ابد ثبت کرد (حافظه متناهی)
\item
  نمی‌توان هر بسته را اندازه‌گیری کرد (تعداد بسیار زیاد)
\end{itemize}

هدف: \textbf{استنباط ویژگی‌های جامعه از نمونه.}

\PSsection{12.3}{پارامتر در برابر آماره}

\PSsubsection{پارامتر}

یک عدد \textbf{ثابت} (اما نامعلوم) که جامعه را توصیف می‌کند:

\begin{itemize}
\tightlist
\item
  \lr{$\mu$} = میانگین واقعی جامعه
\item
  \lr{$\sigma^{2}$} = واریانس واقعی جامعه
\item
  \lr{$p$} = نرخ واقعی نقص
\item
  \lr{$\lambda$} = نرخ واقعی خرابی
\end{itemize}

پارامترها \textbf{ثابت} هستند \lr{—} یک مقدار واقعی دارند (فقط آن را نمی‌دانیم).

\PSsubsection{آماره}

یک \textbf{تابع از داده نمونه} که برای برآورد یا آزمون پارامترها به کار می‌رود:

\begin{itemize}
\tightlist
\item
  \lr{$\bar{X}$} = میانگین نمونه (برآورد \lr{$\mu$})
\item
  \lr{$S^{2}$} = واریانس نمونه (برآورد \lr{$\sigma^{2}$})
\item
  \lr{$\hat{p}$} = نسبت نمونه (برآورد \lr{$p$})
\end{itemize}

\begin{PSquote}

\lr{\mbox{\PSwarn{}}} \textbf{نکته کلیدی:} آماره‌ها متغیرهای تصادفی هستند! نمونه‌های مختلف مقادیر مختلفی می‌دهند.

\end{PSquote}

\begin{PSbox}[unbreakable]{ex}{مثال}

توان واقعی نویز: \lr{$\sigma^{2} = 0.04\,\mathrm{V}^{2}$} (پارامتر \lr{—} ثابت، نامعلوم)

آزمایش 1: جمع‌آوری 100 نمونه \lr{\mbox{\ensuremath{\to}}} محاسبه \lr{$S^{2} = 0.038$} (یک تحقق)\PSbr{}آزمایش 2: جمع‌آوری 100 نمونه \lr{\mbox{\ensuremath{\to}}} محاسبه \lr{$S^{2} = 0.042$} (تحقق متفاوت)\PSbr{}آزمایش 3: جمع‌آوری 100 نمونه \lr{\mbox{\ensuremath{\to}}} محاسبه \lr{$S^{2} = 0.041$} (تحقق دیگر)

هر \lr{$S^{2}$} به‌دلیل تصادفی‌بودن نمونه‌گیری متفاوت است.

\end{PSbox}

\PSsection{12.4}{نمونه تصادفی}

\begin{PSbox}[unbreakable]{def}{تعریف}

یک \textbf{نمونه تصادفی} \lr{$X_{1},\, X_{2},\, \ldots,\, X_{n}$} مجموعه‌ای از \lr{$n$} متغیر تصادفی مستقل و هم‌توزیع \lr{$(i.i.d.)$} از توزیع جامعه است.

\end{PSbox}

\PSsubsection{شرایط لازم}

\begin{enumerate}
\def\labelenumi{\arabic{enumi}.}
\tightlist
\item
  \textbf{مستقل:} هر مشاهده بر سایرین اثر نمی‌گذارد
\item
  \textbf{هم‌توزیع:} همه از همان جامعه آمده‌اند
\item
  \textbf{انتخاب تصادفی:} بدون بایاس نظام‌مند در انتخاب اقلام مورد اندازه‌گیری
\end{enumerate}

\PSsubsection{زمانی که \lr{\mbox{i.i.d}}. برقرار است}

\lr{\mbox{\PSyes{}}} انتخاب تصادفی قطعات از یک دسته بزرگ\PSbr{}\lr{\mbox{\PSyes{}}} اندازه‌گیری‌های مستقل از یک فرایند پایدار\PSbr{}\lr{\mbox{\PSyes{}}} نمونه‌های نویز با فاصله بیش از زمان همبستگی

\PSsubsection{زمانی که \lr{\mbox{i.i.d}}. برقرار نیست}

\lr{\mbox{\PSno{}}} نمونه‌های متوالی یک سیگنال همبسته\PSbr{}\lr{\mbox{\PSno{}}} اندازه‌گیری‌ها در طول یک فرایند رانش‌دار\PSbr{}\lr{\mbox{\PSno{}}} انتخاب غیرتصادفی (مثلاً آزمایش فقط قطعات «مشکوک»)

\PSsection{12.5}{توزیع‌های نمونه‌گیری}

\begin{PSbox}[unbreakable]{def}{تعریف}

\textbf{توزیع نمونه‌گیری} یک آماره همان توزیع احتمال آن آماره در تمام نمونه‌های ممکن به اندازه \lr{$n$} است.

\end{PSbox}

\PSsubsection{چرا اهمیت دارد}

یک آماره (مانند \lr{$\bar{X}$}) یک متغیر تصادفی است. توزیع آن به ما می‌گوید:

\begin{itemize}
\tightlist
\item
  چقدر از نمونه‌ای به نمونه دیگر تغییر می‌کند؟
\item
  با چه احتمالی به پارامتر واقعی نزدیک است؟
\item
  عدم‌قطعیت خود را چگونه کمّی کنیم؟
\end{itemize}

این بنیاد بازه‌های اطمینان و آزمون فرض است.

\PSsection{12.6}{میانگین نمونه \lr{$\bar{X}$}}

\begin{PSbox}[unbreakable]{def}{تعریف}

\PSdisplay{\bar{X} = \frac{1}{n}\sum_{i=1}^n X_i}

\end{PSbox}

\PSsubsection{ویژگی‌ها}

\textbf{میانگین \lr{$\bar{X}$}:}\PSdisplay{E[\bar{X}] = \mu}

\lr{$\bar{X}$} یک برآوردگر \textbf{نااریب} برای \lr{$\mu$} است (به‌طور میانگین برابر مقدار واقعی است).

\textbf{واریانس \lr{$\bar{X}$}:}\PSdisplay{\PSmtext{Var}(\bar{X}) = \frac{\sigma^2}{n}}

اندازه‌گیری‌های بیشتر \lr{\mbox{\ensuremath{\to}}} تغییرپذیری کمتر \lr{\mbox{\ensuremath{\to}}} برآورد دقیق‌تر.

\textbf{خطای معیار:}\PSdisplay{\PSmtext{SE}(\bar{X}) = \frac{\sigma}{\sqrt{n}}}

\PSsubsection{توزیع نمونه‌گیری \lr{$\bar{X}$}}

\begin{itemize}
\tightlist
\item
  اگر جامعه گاوسی باشد: \lr{$\bar{X} \sim N\left(\mu,\, \frac{\sigma^{2}}{n}\right)$} \textbf{به‌طور دقیق} برای هر \lr{$n$}
\item
  اگر جامعه گاوسی نباشد: \lr{$\bar{X} \approx N\left(\mu,\, \frac{\sigma^{2}}{n}\right)$} \textbf{به‌طور تقریبی} برای \lr{$n$} بزرگ (بر اساس \lr{\mbox{CLT}})
\end{itemize}

\PSsection{12.7}{واریانس نمونه \lr{$S^{2}$}}

\begin{PSbox}[unbreakable]{def}{تعریف}

\PSdisplay{S^2 = \frac{1}{n-1}\sum_{i=1}^n (X_i - \bar{X})^2}

\end{PSbox}

\PSsubsection{چرا \lr{$n-1$}؟ (تصحیح بسل)}

استفاده از \lr{$n$} در مخرج یک برآوردگر \textbf{اریب} می‌دهد (به‌طور نظام‌مند \lr{$\sigma^{2}$} را کم‌برآورد می‌کند):

\lr{$\displaystyle E\left[\frac{1}{n} \cdot \sum (X_{i} - \bar{X})^{2}\right] = \frac{n-1}{n}\sigma^{2} \ne \sigma^{2}$}

استفاده از \lr{$n-1$} این را تصحیح می‌کند:\PSdisplay{E[S^2] = \sigma^2}

\textbf{شهود:} \lr{$\bar{X}$} از همان داده محاسبه می‌شود، بنابراین انحراف‌ها از \lr{$\bar{X}$} کوچک‌تر از انحراف‌ها از \lr{$\mu$} هستند. با برآورد میانگین، «یک درجه آزادی را از دست می‌دهیم».

\PSsubsection{درجات آزادی}

با \lr{$n$} نقطه داده و برآورد \lr{$\bar{X}$}:

\begin{itemize}
\tightlist
\item
  \lr{$n$} انحراف \lr{$(X_{i} - \bar{X})$} جمعاً صفر می‌شوند: \lr{$\sum (X_{i} - \bar{X}) = 0$}
\item
  تنها \lr{$n-1$} تای آن‌ها «آزاد» هستند \lr{—} آخری تعیین‌شده است
\item
  \lr{$df = n - 1$}
\end{itemize}

\PSsubsection{توزیع نمونه‌گیری (جامعه نرمال)}

اگر جامعه \lr{$N(\mu,\, \sigma^{2})$} باشد:\PSdisplay{(n-1)S^2/\sigma^2 \sim \chi^2(n-1)}

\PSsection{12.8}{انحراف معیار نمونه}

\PSdisplay{S = \sqrt{S^2} = \sqrt{\frac{1}{n-1}\sum_{i=1}^n (X_i - \bar{X})^2}}

توجه: \lr{$E[S] \ne \sigma$} (\lr{\mbox{S}} برای \lr{$\sigma$} کمی اریب است، هرچند \lr{$S^{2}$} برای \lr{$\sigma^{2}$} نااریب است).

\PSsection{12.9}{ارتباط با \lr{\mbox{CLT}}}

از ماژول 11: \lr{$\bar{X} \to N\left(\mu,\, \frac{\sigma^{2}}{n}\right)$} وقتی \lr{$n \to \infty$}.

این بدان معناست که توزیع نمونه‌گیری \lr{$\bar{X}$} تقریباً گاوسی است. این همان چیزی است که به ما امکان می‌دهد:

\begin{enumerate}
\def\labelenumi{\arabic{enumi}.}
\tightlist
\item
  بازه‌های اطمینان بسازیم (ماژول 14)
\item
  آزمون‌های فرض انجام دهیم (ماژول 15)
\item
  عدم‌قطعیت در برآوردها را کمّی کنیم
\end{enumerate}

\PSsection{12.10}{مثال‌های \lr{\mbox{MATLAB}}}

\begin{PSbox}[unbreakable]{ex}{مثال 1: توزیع نمونه‌گیری \lr{$\bar{X}$}}

\tcbset{PScodemode/.style={unbreakable}}

\begin{Shaded}
\begin{Highlighting}[]
\CommentTok{\%\% Demonstrate: X̄ is a random variable with its own distribution}
\CommentTok{\% Population: Exponential(λ=1), μ=1, σ²=1}
\VariableTok{mu\_pop} \OperatorTok{=} \FloatTok{1}\OperatorTok{;} \VariableTok{sigma2\_pop} \OperatorTok{=} \FloatTok{1}\OperatorTok{;}
\VariableTok{n} \OperatorTok{=} \FloatTok{25}\OperatorTok{;}  \CommentTok{\% Sample size}
\VariableTok{N\_experiments} \OperatorTok{=} \FloatTok{10000}\OperatorTok{;}

\CommentTok{\% Repeat experiment N\_experiments times}
\VariableTok{X\_bar} \OperatorTok{=} \VariableTok{zeros}\NormalTok{(}\FloatTok{1}\OperatorTok{,} \VariableTok{N\_experiments}\NormalTok{)}\OperatorTok{;}
\KeywordTok{for} \VariableTok{i} \OperatorTok{=} \FloatTok{1}\OperatorTok{:}\VariableTok{N\_experiments}
    \VariableTok{sample} \OperatorTok{=} \VariableTok{exprnd}\NormalTok{(}\FloatTok{1}\OperatorTok{,} \FloatTok{1}\OperatorTok{,} \VariableTok{n}\NormalTok{)}\OperatorTok{;}
    \VariableTok{X\_bar}\NormalTok{(}\VariableTok{i}\NormalTok{) }\OperatorTok{=} \VariableTok{mean}\NormalTok{(}\VariableTok{sample}\NormalTok{)}\OperatorTok{;}
\KeywordTok{end}

\VariableTok{figure}\OperatorTok{;}
\VariableTok{histogram}\NormalTok{(}\VariableTok{X\_bar}\OperatorTok{,} \FloatTok{60}\OperatorTok{,} \SpecialStringTok{\textquotesingle{}Normalization\textquotesingle{}}\OperatorTok{,} \SpecialStringTok{\textquotesingle{}pdf\textquotesingle{}}\NormalTok{)}\OperatorTok{;}
\VariableTok{hold} \VariableTok{on}\OperatorTok{;}
\VariableTok{x} \OperatorTok{=} \VariableTok{linspace}\NormalTok{(}\FloatTok{0}\OperatorTok{,} \FloatTok{3}\OperatorTok{,} \FloatTok{200}\NormalTok{)}\OperatorTok{;}
\VariableTok{plot}\NormalTok{(}\VariableTok{x}\OperatorTok{,} \VariableTok{normpdf}\NormalTok{(}\VariableTok{x}\OperatorTok{,} \VariableTok{mu\_pop}\OperatorTok{,} \VariableTok{sqrt}\NormalTok{(}\VariableTok{sigma2\_pop}\OperatorTok{/}\VariableTok{n}\NormalTok{))}\OperatorTok{,} \SpecialStringTok{\textquotesingle{}r\textquotesingle{}}\OperatorTok{,} \SpecialStringTok{\textquotesingle{}LineWidth\textquotesingle{}}\OperatorTok{,} \FloatTok{2}\NormalTok{)}\OperatorTok{;}
\VariableTok{xlabel}\NormalTok{(}\SpecialStringTok{\textquotesingle{}Sample Mean\textquotesingle{}}\NormalTok{)}\OperatorTok{;} \VariableTok{ylabel}\NormalTok{(}\SpecialStringTok{\textquotesingle{}PDF\textquotesingle{}}\NormalTok{)}\OperatorTok{;}
\VariableTok{title}\NormalTok{(}\VariableTok{sprintf}\NormalTok{(}\SpecialStringTok{\textquotesingle{}Sampling Distribution of X̄ (n=\%d)\textquotesingle{}}\OperatorTok{,} \VariableTok{n}\NormalTok{))}\OperatorTok{;}
\VariableTok{legend}\NormalTok{(}\SpecialStringTok{\textquotesingle{}Empirical\textquotesingle{}}\OperatorTok{,} \VariableTok{sprintf}\NormalTok{(}\SpecialStringTok{\textquotesingle{}N(\%.1f, \%.4f)\textquotesingle{}}\OperatorTok{,} \VariableTok{mu\_pop}\OperatorTok{,} \VariableTok{sigma2\_pop}\OperatorTok{/}\VariableTok{n}\NormalTok{))}\OperatorTok{;}

\VariableTok{fprintf}\NormalTok{(}\SpecialStringTok{\textquotesingle{}E[X̄] = \%.4f (theory: \%.4f)\textbackslash{}n\textquotesingle{}}\OperatorTok{,} \VariableTok{mean}\NormalTok{(}\VariableTok{X\_bar}\NormalTok{)}\OperatorTok{,} \VariableTok{mu\_pop}\NormalTok{)}\OperatorTok{;}
\VariableTok{fprintf}\NormalTok{(}\SpecialStringTok{\textquotesingle{}Var(X̄) = \%.4f (theory: \%.4f)\textbackslash{}n\textquotesingle{}}\OperatorTok{,} \VariableTok{var}\NormalTok{(}\VariableTok{X\_bar}\NormalTok{)}\OperatorTok{,} \VariableTok{sigma2\_pop}\OperatorTok{/}\VariableTok{n}\NormalTok{)}\OperatorTok{;}
\end{Highlighting}
\end{Shaded}

\end{PSbox}

\begin{PSbox}[unbreakable]{ex}{مثال 2: توزیع واریانس نمونه}

\tcbset{PScodemode/.style={unbreakable}}

\begin{Shaded}
\begin{Highlighting}[]
\CommentTok{\%\% Sampling Distribution of S²}
\CommentTok{\% Population: N(0, 4), σ² = 4}
\VariableTok{sigma2} \OperatorTok{=} \FloatTok{4}\OperatorTok{;} \VariableTok{n} \OperatorTok{=} \FloatTok{10}\OperatorTok{;}
\VariableTok{N\_experiments} \OperatorTok{=} \FloatTok{10000}\OperatorTok{;}

\VariableTok{S2} \OperatorTok{=} \VariableTok{zeros}\NormalTok{(}\FloatTok{1}\OperatorTok{,} \VariableTok{N\_experiments}\NormalTok{)}\OperatorTok{;}
\KeywordTok{for} \VariableTok{i} \OperatorTok{=} \FloatTok{1}\OperatorTok{:}\VariableTok{N\_experiments}
    \VariableTok{sample} \OperatorTok{=} \VariableTok{sqrt}\NormalTok{(}\VariableTok{sigma2}\NormalTok{) }\OperatorTok{*} \VariableTok{randn}\NormalTok{(}\FloatTok{1}\OperatorTok{,} \VariableTok{n}\NormalTok{)}\OperatorTok{;}
    \VariableTok{S2}\NormalTok{(}\VariableTok{i}\NormalTok{) }\OperatorTok{=} \VariableTok{var}\NormalTok{(}\VariableTok{sample}\NormalTok{)}\OperatorTok{;}   \CommentTok{\% Uses n{-}1 denominator}
\KeywordTok{end}

\VariableTok{fprintf}\NormalTok{(}\SpecialStringTok{\textquotesingle{}E[S²] = \%.4f (theory: \%.4f)\textbackslash{}n\textquotesingle{}}\OperatorTok{,} \VariableTok{mean}\NormalTok{(}\VariableTok{S2}\NormalTok{)}\OperatorTok{,} \VariableTok{sigma2}\NormalTok{)}\OperatorTok{;}
\VariableTok{fprintf}\NormalTok{(}\SpecialStringTok{\textquotesingle{}Var(S²) = \%.4f (theory: \%.4f)\textbackslash{}n\textquotesingle{}}\OperatorTok{,} \VariableTok{var}\NormalTok{(}\VariableTok{S2}\NormalTok{)}\OperatorTok{,} \FloatTok{2}\OperatorTok{*}\VariableTok{sigma2}\OperatorTok{\^{}}\FloatTok{2}\OperatorTok{/}\NormalTok{(}\VariableTok{n}\OperatorTok{{-}}\FloatTok{1}\NormalTok{))}\OperatorTok{;}

\VariableTok{figure}\OperatorTok{;}
\VariableTok{histogram}\NormalTok{(}\VariableTok{S2}\OperatorTok{,} \FloatTok{80}\OperatorTok{,} \SpecialStringTok{\textquotesingle{}Normalization\textquotesingle{}}\OperatorTok{,} \SpecialStringTok{\textquotesingle{}pdf\textquotesingle{}}\NormalTok{)}\OperatorTok{;}
\VariableTok{hold} \VariableTok{on}\OperatorTok{;}
\VariableTok{x} \OperatorTok{=} \VariableTok{linspace}\NormalTok{(}\FloatTok{0}\OperatorTok{,} \FloatTok{15}\OperatorTok{,} \FloatTok{200}\NormalTok{)}\OperatorTok{;}
\CommentTok{\% (n{-}1)S²/σ² \textasciitilde{} χ²(n{-}1), so S² \textasciitilde{} σ²/(n{-}1) * χ²(n{-}1)}
\VariableTok{plot}\NormalTok{(}\VariableTok{x}\OperatorTok{,} \VariableTok{chi2pdf}\NormalTok{(}\VariableTok{x}\OperatorTok{*}\NormalTok{(}\VariableTok{n}\OperatorTok{{-}}\FloatTok{1}\NormalTok{)}\OperatorTok{/}\VariableTok{sigma2}\OperatorTok{,} \VariableTok{n}\OperatorTok{{-}}\FloatTok{1}\NormalTok{)}\OperatorTok{*}\NormalTok{(}\VariableTok{n}\OperatorTok{{-}}\FloatTok{1}\NormalTok{)}\OperatorTok{/}\VariableTok{sigma2}\OperatorTok{,} \SpecialStringTok{\textquotesingle{}r\textquotesingle{}}\OperatorTok{,} \SpecialStringTok{\textquotesingle{}LineWidth\textquotesingle{}}\OperatorTok{,} \FloatTok{2}\NormalTok{)}\OperatorTok{;}
\VariableTok{xlabel}\NormalTok{(}\SpecialStringTok{\textquotesingle{}S²\textquotesingle{}}\NormalTok{)}\OperatorTok{;} \VariableTok{ylabel}\NormalTok{(}\SpecialStringTok{\textquotesingle{}PDF\textquotesingle{}}\NormalTok{)}\OperatorTok{;}
\VariableTok{title}\NormalTok{(}\VariableTok{sprintf}\NormalTok{(}\SpecialStringTok{\textquotesingle{}Sampling Distribution of S² (n=\%d, σ²=\%d)\textquotesingle{}}\OperatorTok{,} \VariableTok{n}\OperatorTok{,} \VariableTok{sigma2}\NormalTok{))}\OperatorTok{;}
\VariableTok{legend}\NormalTok{(}\SpecialStringTok{\textquotesingle{}Empirical\textquotesingle{}}\OperatorTok{,} \SpecialStringTok{\textquotesingle{}Theoretical (scaled χ²)\textquotesingle{}}\NormalTok{)}\OperatorTok{;}
\end{Highlighting}
\end{Shaded}

\end{PSbox}

\begin{PSbox}[unbreakable]{ex}{مثال 3: اثر اندازه نمونه}

\tcbset{PScodemode/.style={unbreakable}}

\begin{Shaded}
\begin{Highlighting}[]
\CommentTok{\%\% Precision improves with n: Var(X̄) = σ²/n}
\VariableTok{sigma} \OperatorTok{=} \FloatTok{2}\OperatorTok{;}
\VariableTok{n\_values} \OperatorTok{=}\NormalTok{ [}\FloatTok{5}\OperatorTok{,} \FloatTok{10}\OperatorTok{,} \FloatTok{25}\OperatorTok{,} \FloatTok{50}\OperatorTok{,} \FloatTok{100}\OperatorTok{,} \FloatTok{500}\NormalTok{]}\OperatorTok{;}
\VariableTok{N\_exp} \OperatorTok{=} \FloatTok{10000}\OperatorTok{;}

\VariableTok{fprintf}\NormalTok{(}\SpecialStringTok{\textquotesingle{}n\textbackslash{}t| SE(X̄) Theory\textbackslash{}t| SE(X̄) Sim\textbackslash{}n\textquotesingle{}}\NormalTok{)}\OperatorTok{;}
\VariableTok{fprintf}\NormalTok{(}\SpecialStringTok{\textquotesingle{}{-}{-}{-}{-}{-}{-}{-}{-}+{-}{-}{-}{-}{-}{-}{-}{-}{-}{-}{-}{-}{-}{-}{-}+{-}{-}{-}{-}{-}{-}{-}{-}{-}{-}{-}\textbackslash{}n\textquotesingle{}}\NormalTok{)}\OperatorTok{;}
\KeywordTok{for} \VariableTok{n} \OperatorTok{=} \VariableTok{n\_values}
    \VariableTok{X\_bar} \OperatorTok{=} \VariableTok{mean}\NormalTok{(}\VariableTok{sigma}\OperatorTok{*}\VariableTok{randn}\NormalTok{(}\VariableTok{n}\OperatorTok{,} \VariableTok{N\_exp}\NormalTok{)}\OperatorTok{,} \FloatTok{1}\NormalTok{)}\OperatorTok{;}
    \VariableTok{fprintf}\NormalTok{(}\SpecialStringTok{\textquotesingle{}\%d\textbackslash{}t| \%.4f\textbackslash{}t| \%.4f\textbackslash{}n\textquotesingle{}}\OperatorTok{,} \VariableTok{n}\OperatorTok{,} \VariableTok{sigma}\OperatorTok{/}\VariableTok{sqrt}\NormalTok{(}\VariableTok{n}\NormalTok{)}\OperatorTok{,} \VariableTok{std}\NormalTok{(}\VariableTok{X\_bar}\NormalTok{))}\OperatorTok{;}
\KeywordTok{end}
\end{Highlighting}
\end{Shaded}

\end{PSbox}

\PSsection{12.11}{مسایل تمرینی}

\begin{PSbox}[unbreakable]{prob}{مسیله 1}

یک جامعه دارای \lr{$\mu = 50$} و \lr{$\sigma = 10$} است. یک نمونه تصادفی به اندازه \lr{$n = 64$} گرفته می‌شود.

\begin{enumerate}
\def\labelenumi{(\PSalph{enumi})}
\tightlist
\item
  \lr{$E[\bar{X}]$} چقدر است؟ (ب) \lr{$\operatorname{Var}(\bar{X})$} چقدر است؟ (ج) \lr{$P(48 < \bar{X} < 52)$} چقدر است؟
\end{enumerate}

\PSsolution{} 

\begin{enumerate}
\def\labelenumi{(\PSalph{enumi})}
\tightlist
\item
  \lr{$E[\bar{X}] = 50$}
\item
  \lr{$\operatorname{Var}(\bar{X}) = \frac{100}{64} = 1.5625$}، \lr{$\operatorname{SE} = 1.25$}
\item
  \lr{$P(48 < \bar{X} < 52) = P(-1.6 < Z < 1.6) = 2\Phi (1.6) - 1 = 0.8904$}
\end{enumerate}

\end{PSbox}

\begin{PSbox}[unbreakable]{prob}{مسیله 2}

چرا در فرمول واریانس نمونه از \lr{$n-1$} استفاده می‌کنیم و نه \lr{$n$}؟ هم به‌صورت ریاضی و هم به‌صورت شهودی توضیح دهید.

\PSsolution{} از نظر ریاضی: \lr{$E\left[\frac{1}{n}\sum (X_{i}-\bar{X})^{2}\right] = \frac{n-1}{n} \cdot \sigma^{2}$}، بنابراین تقسیم بر \lr{$n-1$} نتیجه می‌دهد \lr{$E[S^{2}] = \sigma^{2}$}. از نظر شهودی: ما از \lr{$\bar{X}$} همان داده استفاده می‌کنیم (نه \lr{$\mu$} واقعی)، بنابراین انحراف‌ها از \lr{$\bar{X}$} به‌طور مصنوعی کوچک‌اند. یک درجه آزادی «صرف» برآورد میانگین می‌شود.

\end{PSbox}

\begin{PSbox}[unbreakable]{prob}{مسیله 3}

یک مهندس توان نویز را از 20 نمونه اندازه‌گیری می‌کند و \lr{$S^{2} = 0.045\,\mathrm{V}^{2}$} به دست می‌آید. واریانس واقعی \lr{$\sigma^{2} = 0.04\,\mathrm{V}^{2}$} است. آیا این نگران‌کننده است؟ برای ارزیابی از توزیع نمونه‌گیری استفاده کنید.

\PSsolution{} تحت \lr{$H_{0}: \sigma^{2} = 0.04$}، داریم \lr{$\frac{(n-1)S^{2}}{\sigma^{2}} \sim \chi^{2}(19)$}. مقدار مشاهده‌شده: \lr{$\frac{19(0.045)}{0.04} = 21.375$. $P(\chi^{2}(19) > 21.375) \approx 0.31$}. این غیرمعمول نیست \lr{—} حدود 31\% از نمونه‌ها هنگامی که \lr{$\sigma^{2} = 0.04$} است، \lr{$S^{2} \ge 0.045$} می‌دهند. نگران‌کننده نیست.

\emph{ماژول بعدی: {ماژول 13 \lr{—} برآورد پارامتر}}

\end{PSbox}
